\section{Related Work}

\subsection{Multimodal Affective Understanding and Conflict}

Multimodal affect datasets such as CH-SIMS v2 provide the empirical basis for studying
how visual, textual, and acoustic cues jointly express sentiment \cite{liu2022chsims}.
CA-MER studies realistic audio--visual emotion conflicts and documents systematic
modality bias in multimodal emotion reasoning \cite{han2025benchmarking}. EmoMM
extends conflict-centered evaluation to aligned, conflict, and missing-modality
settings and analyzes model behavior under these conditions \cite{sun2026emomm}.
Most task formulations ultimately score an output label or description. Our focus is
instead the relation among two unimodal states and their joint state at \(t_0\) under
controlled Conflict and Aligned inputs. Conflict/Aligned is used as input-side
relation supervision; subsequent Misread behavior remains a separate outcome.

\subsection{Pre-generation State Analysis and Uncertainty}

Semantic uncertainty measures invariance in generated meanings
\cite{kuhn2024semanticuncertainty}. In contrast, our state is measured at \(t_0\) from
registered hidden-state trajectories, before any diagnostic description is generated.
HALP shows that pre-generation internal representations can support probes of later
hallucination outcomes \cite{kogilathota-etal-2026-halp}. This supports the broader
premise that subsequent behavior can be readable before decoding, but hallucination
is not affective Misread and our objective is not a detector. We characterize the
three-condition structure first and examine its behavioral association separately.
We do not turn the resulting coordinates into one uncertainty or Misread score.
State Dispersion, Modality Split, and signed Joint Lean describe prompt stability,
unimodal separation, and joint direction, respectively.

\subsection{Frozen Representation Evaluation}

Our learned representations are evaluated only after the Conflict/Aligned encoders
are frozen. A controlled MLP probe is trained and tested within Conflict samples to
measure association with later Misread labels. Single-Point, Trajectory MLP, and TME
share data, supervision budgets, probes, and test intersections, while retaining
their native representation structures and objectives. This evaluates each complete
method rather than attributing a difference to one unablated component.

\subsection{State-guided Response}

State-guided response is an application of state analysis rather than the central
claim. It is evaluated by response-level Affective Misattribution and cue preservation;
it is not presented as correcting the base model's earlier Diagnostic Misread.
